\chapter[Формирования и прием OFDM. Построение спектра радиосигнала (преобразование Фурье). Изучение работы библиотеки FFTW3]{Формирования и прием OFDM. Построение спектра радиосигнала (преобразование Фурье). Изучение работы библиотеки FFTW3.}

\begin{leftbar}
    \href{https://github.com/someotb/sdr}{Ссылка на GitHub}
\end{leftbar}

\section*{Цель практики:}
Сформировать, отправить и принять OFDM символы. Построить спектр радиосигнала. Использовать библиотеку fftw3 для FFT.

\section*{Краткие теоретические сведения}

\textbf{Формирование буффера}\newline
\begin{figure}[H]
    \centering
    \includegraphics[width=1.0\textwidth]{ofdm_symbol.png}
    \caption{Как выгядит OFDM символ в частотной области}
\end{figure}

Такими OFDM символами мы заполняем буффер размером 1920 символов. Туда влезает 1920 / 128 = 15 OFDM символов, но не забываем что у нас еще есть PSS Zadoff-Chu, который
занимает 128 символов, то есть под OFDM символы у нас уходит 1920 - 128 символов. Но, также у нас есть CP - cyclic prefex размером 32 символа, значит размер передоваемого OFDM символа будет 128 + 32 = 160, отсюда общее кол-во
OFDM символов котрое мы сможем передать это 1920 / 160 = 12 - 1(под Zadoff-Chu)

\textbf{Построение спектра. Библиотека fftw3}\newline

Чтобы построить спектр радиосигнала нужно провести прямое преобразовние фурье над буффером, чтобы перевести из временной области в частотную. В этом
мне поможет библиотека fftw3, которая оптимизирована под работу алгоритма Фурье.

\section*{Выполнение}

\textbf{Класс для работы с fftw3}\newline

Я уже показывал как формируется OFDM-символ в предыдущем отчете, так что сильно тут останавливаться не будем, перейдем к спектру сигнала.
Для удобства я создал отдельный класс для работы с fftw3:

\begin{lstlisting}[language=C++,caption=Класс для работы с fftw3]
class FFT_Context
{
    private:
        fftwf_plan plan_forward;
        fftwf_plan plan_backward;

    public:
        int N;
        fftwf_complex *in;
        fftwf_complex *out;
        FFT_Context(int n);
        ~FFT_Context();
        void fft();
        void ifft();
};

FFT_Context::FFT_Context(int n) : N(n)
{
    in = fftwf_alloc_complex(N);
    out = fftwf_alloc_complex(N);

    plan_forward = fftwf_plan_dft_1d(N, in, out, FFTW_FORWARD, FFTW_MEASURE);
    plan_backward = fftwf_plan_dft_1d(N, in, out, FFTW_BACKWARD, FFTW_MEASURE);
}

FFT_Context::~FFT_Context()
{
    fftwf_destroy_plan(plan_forward);
    fftwf_destroy_plan(plan_backward);
    fftwf_free(in);
    fftwf_free(out);
}

void FFT_Context::fft()
{
    fftwf_execute(plan_forward);
}

void FFT_Context::ifft()
{
    fftwf_execute(plan_backward);

    float scale = 1.0f / N;
    for (int i = 0; i < N; ++i)
    {
        out[i][0] *= scale;
        out[i][1] *= scale;
    }
}
\end{lstlisting}

Объект класса создается с входным параметром, размер ожидаемого FFT, есть функции для прямого и обратного преобразования, внутренние массивы fftw3
создаются и удаляются автоматически через конструктор и деструктор. Есть встроенный scale после обратного преобразования Фурье.

\vspace{20pt}

\textbf{Спектр радиосигнала}\newline

Спектр делается с помощью той же библиотеки fftw3

\begin{lstlisting}[language=C++,caption=Функция для построения спектра]
void spectrum(std::vector<std::complex<float>> &in_signal, std::vector<float> &shifted_magnitude, std::vector<float> &argument, FFT_Context &context)
{
    std::vector<float> magnitude(in_signal.size(), 0);

    for (size_t i = 0; i < in_signal.size(); ++i)
    {
        context.in[i][0] = std::real(in_signal[i]) / 32768.0;
        context.in[i][1] = std::imag(in_signal[i]) / 32768.0;
    }

    context.fft();

    for (size_t i = 0; i < in_signal.size(); ++i)
    {
        float real = context.out[i][0];
        float imag = context.out[i][1];
        magnitude[i] = 20.0 * log10(sqrt(real * real + imag * imag) / in_signal.size());
        argument[i] = atan2(imag, real);
    }

    for (size_t i = 0; i < in_signal.size() / 2; ++i)
    {
        shifted_magnitude[i] = magnitude[i + in_signal.size() / 2];
        shifted_magnitude[i + in_signal.size() / 2] = magnitude[i];
    }
}
\end{lstlisting}

Функция выполняет прямое преобразование Фурье, затем берет амплитуду сигнала и переводит ее в dB, также она берет phase. Функция
также сдвигает частоты в более привычное расположение, так как после FFT частоты расположены так: [0, +freq, -freq], функция делает их такими:
[-freq, 0, +freq]. Вот как это выглядит в GUI, в правом нижнем углу спектр сигнала QPSK:

\begin{figure}[H]
    \centering
    \includegraphics[width=1.0\textwidth]{spectre.jpg}
    \caption{Спектр сигнала}
\end{figure}

\section*{Вывод}
Был написал класс для удобной работы с fftw3, функция построения спектра, вывод его в GUI.
